\section{Capaian Pembelajaran Lulusan (CPL)}

Capaian Pembelajaran Lulusan (CPL) adalah kompetensi yang harus dimiliki lulusan Program Studi Teknik Informatika. Mata kuliah Pemrograman Web membebankan dua CPL pada mahasiswa, yang mencakup kompetensi dalam aspek pengetahuan, keterampilan, dan sikap. Pemahaman terhadap CPL membantu Anda menempatkan mata kuliah ini dalam konteks yang lebih luas---yaitu kontribusinya terhadap profil lulusan program studi \cite{aptikomObe}.

CPL yang dibebankan pada mata kuliah Pemrograman Web adalah:

\begin{itemize}[leftmargin=*]
  \item \textbf{CPL-04:} Kemampuan mendesain, mengimplementasi dan mengevaluasi solusi berbasis computing yang memenuhi kebutuhan-kebutuhan computing pada sebuah disiplin program.
  \item \textbf{CPL-05:} Mengidentifikasi dan menganalisis kebutuhan-kebutuhan pengguna dan mempertimbangkannya dalam memilih, membuat, mengintegrasi, mengevaluasi, dan mengadministrasi sistem berbasis computing.
\end{itemize}

Implikasi CPL-04 untuk mata kuliah ini ialah bahwa Anda akan belajar mendesain solusi web, mengimplementasikannya dengan HTML, CSS, JavaScript, PHP, dan MySQL, serta mengevaluasi keamanan dan performanya. CPL-05 menuntut Anda mampu mengidentifikasi dan menganalisis kebutuhan pengguna---sesuai dengan Sub-CPMK 1.1---serta mempertimbangkan kebutuhan tersebut dalam memilih teknologi, membuat antarmuka, mengintegrasikan front-end dan back-end, mengevaluasi hasil, dan mengadministrasi aplikasi web.
